\subsection{Angles.}

We shall assume, in this section and the next, that A is an ordered field, hence of characteristic zero. Recall (Rectifications to Chap.~VI) that, if F is a vector space over A, the relation “ there exists λ \textgreater{} 0 such that y = λx ” is an equivalence relation between elements x, y of F - \{0\}, that every equivalence class for this relation is called an open half-line with origin 0, and that the union of an open half-line and \{0\} is called a closed half-line (or simply half-line) with origin 0; if D is a line and Δ a closed half-line contained in D, D is the union of Δ and -Δ, and contains no other closed half-line. We shall say that a half-line is isotropic if the line containing it is isotropic.

Recall also (ibid.) that, given a vector space of finite dimension n over A, an orientation on F is the choice of one of the two half-lines of the space \(\bigwedge^{n}F\) the n-vectors belonging to this half-line are said to be positive. A finite-dimensional vector space equipped with an orientation is said to be oriented.

The homotheties of E whose ratio is \textgreater0 obviously form a subgroup of index 2 in H; we will denote it by H \(^{+}\) . The canonical homomorphism i : O \(^{+}\) → S \(^{+}\) /H (cf.~prop. 2) is the composite of the canonical homomorphisms of S \(^{+}\) /H \(^{+}\) on S \(^{+}\) /H and of O \(^{+}\) in S \(^{+}\) /H \(^{+}\) ; since O \(^{+}\) ∩ H \(^{+}\) = \{1\} this latter homomorphism is injective.

PROPOSITION 7. --- Suppose that A is a maximal ordered field and that \(\Phi\) be a positive form (§ 7). Then the canonical homomorphisms of O \(^{+}\) into S \(^{+}\) /H \(^{+}\) and of O \(^{+}\) /\{1, -1\} into S \(^{+}\) /H are bijective, and S \(^{+}\) is isomorphic to O \(^{+}\) × H \(^{+}\) .

We have already seen that the homomorphisms in question are injective, and it suffices to show that the first one is surjective. Let \((e_{1}, e_{2})\) an orthonormal basis of E, and let w be the direct similarity such that \(w(e_{1}) = e_{2}\) (cor. 1 of prop. 1, no. 1); we then have \(w^{2} = -1\) (prop. 1, b)). Given any direct similarity transformation \(u = a + bw\) ( \(a \in \mathbf{A}, b \in \mathbf{A}\) ), we have \(\mathrm{N}(u) = a^2 + b^2 > 0\) and there exists one and only one rotation contained in the same half-line of \(\mathbf{A}(\Phi)\) as \(u\) , namely \((a^2 + b^2)^{1/2} u\) . QED.

COROLLARY. --- Given two half-lines D, D' with origin 0, there exists one and only one rotation v such that v(D) = D'.

The hypotheses indeed imply that E contains no isotropic lines. Our assertion then follows from cor. 1 of prop. 1 (no. 1).

We shall henceforth assume that A is a maximal ordered field, and that the form \(\Phi\) is positive. In the set of pairs \((\mathrm{D}_1, \mathrm{D}_2)\) of lines (resp. half-lines with origin 0) of E, the relation “there exists a direct similarity (resp. a rotation) u such that \(u(\mathrm{D}_1) = \mathrm{D}_1'\) and \(u(\mathrm{D}_2) = \mathrm{D}_2'\) ” is an equivalence relation between pairs \((\mathrm{D}_1, \mathrm{D}_2)\) and \((\mathrm{D}_1', \mathrm{D}_2')\) The equivalence class of the pair \((\mathrm{D}_1, \mathrm{D}_2)\) is called, by definition, the angle of the lines (resp. half-lines) \(\mathrm{D}_1, \mathrm{D}_2\) (taken in this order); we denote it \((\widehat{\mathrm{D}_1, \mathrm{D}_2})\) .

PROPOSITION 8. --- Suppose that A is a maximal ordered field, and that the form \(\Phi\) is positive. Let \(D_{1}, D_{2}, D_{1}^{\prime}, D_{2}^{\prime}\) be four lines (resp. half-lines) with origin 0 in E. For the angles \((\widehat{D_{1}, D_{2}})\) and \((\widehat{D_{1}^{\prime}, D_{2}^{\prime}})\) to be equal, it is necessary and sufficient that the angles \((\widehat{D_{1}, D_{1}^{\prime}})\) and \((\widehat{D_{2}, D_{2}^{\prime}})\) be equal.

Let us prove the necessity of the stated condition. Let \(u\) be a direct similarity (resp. a rotation) such that \(u(\mathrm{D}_{1}) = \mathrm{D}_{1}^{\prime}\) and \(u(\mathrm{D}_{2}) = \mathrm{D}_{2}^{\prime}\) . By cor. 3 of prop. 1, there exists a direct similarity (resp. by the corollary of prop. 7, a rotation) \(\nu\) such that \(\nu(\mathrm{D}_{1}) = \mathrm{D}_{2}\) . Since the group \(\mathrm{S}^{+}\) (resp. \(\mathrm{O}^{+}\) ) is commutative, we have \(\mathrm{D}_{2}^{\prime} = u(\nu(\mathrm{D}_{1})) = \nu(u(\mathrm{D}_{1})) = \nu(\mathrm{D}_{1}^{\prime})\) , and this shows that \((\widehat{\mathrm{D}_{1}, \mathrm{D}_{1}^{\prime})} = (\widehat{\mathrm{D}_{2}, \mathrm{D}_{2}^{\prime}})\) . Sufficiency follows from necessity by interchanging \(\mathrm{D}_{2}\) and \(\mathrm{D}_{1}^{\prime}\) .

It follows from prop. 8 that, to every angle \((\widehat{\mathrm{D}_1,\mathrm{D}_2})\) of lines (resp. half-lines) with origin 0 in E, there is canonically associated a well-determined element of \(\mathrm{S}^{+}/\mathrm{H}\) (resp. \(\mathrm{O}^{+}\) ), namely the class mod. H of the direct similarities \(\nu\) (resp. the rotation \(\nu\) ) such that \(u(D_1) = D_2\) for any representative ( \(D_1, D_2\) ) of the angle ( \(\widehat{D_1, D_2}\) ). We have thus defined a canonical bijection \(h\) (resp. \(h'\) ) from the set \(\mathfrak{A}_0\) of angles of lines (resp. \(\mathfrak{A}\) of angles of half-lines) onto \(S^+/H\) (resp. \(O^+\) ); in particular, for every \(\varphi \in \mathfrak{A}\) , one says that \(h(\varphi)\) is the rotation with angle \(\varphi\) . We shall transport to \(\mathfrak{A}_0\) and \(\mathfrak{A}\) , by means of \(h^{-1}\) and of \(h'^{-1}\) , the structures of commutative groups of \(S^+/H\) and of \(O^+\) , and we shall denote additively the groups \(\mathfrak{A}_0\) and \(\mathfrak{A}\) thus obtained. If one denotes by D, \(D', D''\) lines (resp. half-lines) with origin 0 in E, one has by definition

\[
\widehat {(\mathbf {D} , \mathbf {D} ^ {\prime \prime})} = \widehat {(\mathbf {D} , \mathbf {D} ^ {\prime})} + \widehat {(\mathbf {D} ^ {\prime} , \mathbf {D} ^ {\prime \prime})}\tag{9}
\]

(Chasles relation);

one deduces from this

\[
\widehat {(\mathrm{D} , \mathrm{D})} = 0, \quad \widehat {(\mathrm{D} , \mathrm{D} ^ {\prime})} = - (\widehat {\mathrm{D} ^ {\prime} , \mathrm{D}}).\tag{10}
\]

Remarks. --- 1) The set L of lines (resp. half-lines) with origin 0 in E is a homogeneous space of the abelian group S+/H (resp. O+) such that the identity element is the only operator leaving all elements of L invariant. One can therefore apply to L the formulas of chap.~II, 2 \(^{e}\) ed., App. II, n° 1; prop. 8 is thus a particular case of the "parallelogram rule", and formulas (9) and (10) are particular cases of formulas (2) (ibid.).

\begin{enumerate}
\def\labelenumi{\arabic{enumi})}
\setcounter{enumi}{1}
\tightlist
\item
  In the definition of the group of angles of lines, one can, instead of the group S+/H, use the group O+/\{−1, 1\} which is canonically isomorphic to it (prop. 7). The canonical homomorphism of O+ onto O+/\{−1, 1\} thus corresponds to a homomorphism of A onto A0, namely the one which, to the angle of the two half-lines Δ, Δ', associates the angle of the lines D, D' containing respectively Δ, Δ'. *In the case where the field A is the field of real numbers, the group A is thus a double covering of the group A0.*
\end{enumerate}

By prop. 8, all angles of lines (resp. half-lines) of the form \((\widehat{\mathbf{D}^{\prime},\mathbf{D}^{\prime\prime}})\) where \(\mathbf{D}'\) and \(\mathbf{D}''\) are orthogonal (resp.

of the form \((\widehat{\mathrm{D}}, -\mathrm{D})\) are equal: this indeed follows from Remark 2 of No. 1 (resp. is obvious). This angle of lines (resp. of half-lines) is called the right angle (resp. the straight angle); it is an element of order 2 of \(A_{0}\) (resp. of A).

PROPOSITION 9. --- Suppose that the field A is maximally ordered, and that \(\Phi\) is a positive form. For every integer \(n>1\) , the number of elements \(\theta\) of the group \(\mathfrak{A}_{0}\) of angles of lines (resp. \(\mathfrak{A}\) of angles of half-lines) such that \(n\theta=0\) is equal to \(n\) .

As \(\mathfrak{A}_{0}\) , \(\mathrm{S}^{+}/\mathrm{H}\) , \(\mathfrak{A}\) and \(\mathrm{O}^{+}\) are isomorphic (prop. 3), it suffices to carry out the proof for \(\mathrm{O}^{+}\) , that is, to show that there are exactly \(n\) rotations \(\varrho\) such that \(\varphi^{n}=1\) Since A is a maximal ordered field and that \(\mathrm{A}(\Phi)\) is a degree 2 extension field of A (prop. 1 a)), \(\mathrm{A}(\Phi)\) is an algebraically closed field (chap.~VI, § 2, no. 6, th. 3). Hence the roots \(n\) -th roots of unity in \(\mathrm{A}(\Phi)\) form a cyclic group of order \(n\) (chap. V, § 11, no. 1, th. 1). Since we have \(\mathrm{N}(u)=u\overline{u}\geqslant0\) for all \(u\in\mathrm{A}(\Phi)\) , the relation \(u^{n}=1\) implies that one has \(\mathrm{N}(u)=1\) , hence that \(u\) is a rotation (No. 1, Remark 1). This proves our assertion.

COROLLARY. --- The right angle (resp. straight angle) is the only element of order 2 in the group. \(\mathfrak{A}_{0}\) (resp. \(\mathfrak{A}\) ).

We will assume finally that the plane E is oriented.

Lemma 1. --- Let u be a direct similarity of E; all bivectors of the form \(x \wedge u(x)\) belong to the same closed half-line of \(\bigwedge^{2}\mathbf{E}\) .

The case where \(u\) is a homothety is trivial. Otherwise, we have \(x \wedge u(x) \neq 0\) for all \(x \neq 0\) ; let \(x, y\) be two vectors of \(E\) ( \(x \neq 0, y \neq 0\) ); there exists \(v \in S^+\) such that \(y = v(x)\) , whence \(y \wedge u(y) = v(x) \wedge uv(x) = v(x) \wedge v(u(x)) = (\det v)(x \wedge u(x))\) ; by taking an orthonormal basis of \(E\) , we see that det \(v\) is positive (prop. 1 b)); hence our assertion.

That being said, among the two generators \(w\) of \(\mathrm{A}(\Phi)\) such that \(w^{2} = -1\) there exists one and only one such that the bivector \(x \wedge w(x)\) be positive for all \(x\in\mathbf{E}\) . It is this generator that we will choose to define the functions \(c_{w}, s_{w}\) and \(t_{w}\) (no. 2). Let \(h\) and \(h'\) the canonical bijections defined above from the group \(\mathfrak{A}_{0}\) angles of lines on S+/H and of the group \(\mathfrak{A}\) of half-line angles on O+. The composed mappings \(t_{w}\circ h\) of \(\mathfrak{A}_{0}\) in the projective field \(\tilde{\mathbf{A}}\) , \(c_{w}\circ h'\) and \(s_{w}\circ h'\) of \(\mathfrak{A}\) in the field A are denoted respectively by tg, cos, and sin, and are called the tangent, cosine, and sine functions. The mapping \(\varphi\to1/\mathrm{tg}\varphi\) of \(\mathfrak{A}_{0}\) in \(\tilde{\mathbf{A}}\) is denoted cotg and is called the cotangent function. The functions cosine, sine, tangent, and cotangent are said to be the trigonometric functions. The composite mappings tg∘p and cotg∘p, where p denotes the canonical homomorphism from \(\mathfrak{A}\) on \(\mathfrak{A}_{0}\) (Remark 2 above) are still denoted tg and cotg by abuse of language.

Formulas (2), (8), (3), (4), (5), and (7) of §2 give, since here we have \(\delta = -1\)

\begin{enumerate}
\def\labelenumi{(\arabic{enumi})}
\setcounter{enumi}{10}
\tightlist
\item
\end{enumerate}

\[
\operatorname{tg} \left(\varphi + \varphi^ {\prime}\right) = (\operatorname{tg} \varphi + \operatorname{tg} \varphi^ {\prime}) / (1 - \operatorname{tg} \varphi \operatorname{tg} \varphi^ {\prime})\tag{12}
\]

\[
\operatorname{tg} (2 \varphi) = 2 \operatorname{tg} \varphi / (1 - \operatorname{tg} ^ {2} \varphi)
\]

for \(\varphi, \varphi'\) in \(\mathfrak{A}_0\) ;

\begin{enumerate}
\def\labelenumi{(\arabic{enumi})}
\setcounter{enumi}{12}
\tightlist
\item
\end{enumerate}

\[
\cos^ {2} \theta + \sin^ {2} \theta = 1\tag{14}
\]

\[
\cos (\theta + \theta^ {\prime}) = \cos \theta \cos \theta^ {\prime} - \sin \theta \sin \theta^ {\prime}\tag{15}
\]

\[
\sin (\theta + \theta^ {\prime}) = \sin \theta \cos \theta^ {\prime} + \cos \theta \sin \theta^ {\prime}\tag{16}
\]

\[
\sin (2 \theta) = 2 \operatorname{tg} \theta / (1 + \operatorname{tg} ^ {2} \theta),
\]

\[
\cos (2 \theta) = (1 - \mathrm{tg} ^ {2} \theta) / (1 + \mathrm{tg} ^ {2} \theta)
\]

for \(\theta, \theta'\) in \(\mathfrak{A}\) On the other hand, we have, by definition or as an easy consequence of the preceding formulas:

\begin{enumerate}
\def\labelenumi{(\arabic{enumi})}
\setcounter{enumi}{16}
\tightlist
\item
\end{enumerate}

\[
\operatorname{tg} \theta = \sin \theta / \cos \theta , \quad \cot \mathrm{g} \theta = \cos \theta / \sin \theta\tag{18}
\]

\[
1 + \mathrm{tg} ^ {2} \theta = 1 / \cos^ {2} \theta , \quad 1 + \cot \mathrm{g} ^ {2} \theta = 1 / \sin^ {2} \theta
\]

for θ ∈ 𝒜.

Given two nonzero vectors \(x, y\) of E, the angle between these two vectors (taken in this order) is called, and denoted \((\widehat{x}, y)\) the angle of the half-lines to which they belong. For every vector \(x\) from E we call the length of \(x\) is called, and denoted \(|x|\) , the element \(\Phi(x, x)^{1/2}\) of A.

PROPOSITION 10. --- Suppose that the field A is maximally ordered, that the plane E is oriented, and that the form \(\Phi\) is positive. For every pair of nonzero vectors \(x\) , \(y\) of E on \(a\)

\begin{enumerate}
\def\labelenumi{(\arabic{enumi})}
\setcounter{enumi}{18}
\tightlist
\item
\end{enumerate}

\[
\cos \widehat {(x , y)} = \Phi (x, y) / | x |. | y |\tag{20}
\]

\[
\widehat {\sin (x , y)}. e = (x \wedge y) / | x |. | y |,
\]

where e denotes the positive bivector such that \(\Phi_{(2)}(e, e) = 1\) .

Indeed, since the vectors \(x' = x / |x|\) and \(y' = y / |y|\) are both of length 1, there exists a rotation \(\nu\) and exactly one such that \(\nu(x') = y'\) (No. 1, cor. 2 of prop. 1). If we set \(\nu = a + bw\) ( \(a, b\) In A), we have by definition \(a = \cos(\widehat{x,y})\) and \(b = \sin(\widehat{x,y})\) . The relation \(y' = \nu(x') = ax' + bw(x')\) gives \(\Phi(x', y') = a\Phi(x', x') = a\) since \(x'\) and \(w(x')\) are orthogonal (Remark 2 of No. 1); this proves (19). On the other hand, this relation also gives \(x' \wedge y' = bx' \wedge w(x') = b.e\) by the definition of the extension \(\Phi_{(2)}\) of \(\Phi\) à \(\bigwedge^{2} E\) (§ 1, no. 9, formula (37)) and the choice of \(w\) ; this proves (20).

Remarks. --- 3) Given two lines (resp. rays) D, D' of an affine plane L attached to E, one calls the angle of D and D', denoted (D, D'), the angle made by their directions in E (resp. the rays with origin 0 of E corresponding to D and D') (chap.~II, 2). \(^{e}\) ed., App. II, nos. 1 and 3).

\begin{enumerate}
\def\labelenumi{\arabic{enumi})}
\setcounter{enumi}{3}
\tightlist
\item
  Let F be a vector space of arbitrary dimension over the maximal ordered field A, and \(\Psi\) a positive non-degenerate symmetric bilinear form on F. Given two vectors \(x\) , \(y\) linearly independent elements of F, let \(\mathbf{F}'\) the vector plane they span; we call the angle of \(x\) and \(y\) the angle of \(x\) and \(y\) considered as elements of the plane \(\mathbf{F}'\) ; it is denoted \((x,y)\) The cosine of this angle is, by virtue of (19), given by
\end{enumerate}

\[
\cos \widehat {(x , y)} = \Psi (x, y) / | x |. | y |\tag{21}
\]

\((\text{ou } |x| = \Psi(x, x)^{1/2}\) is also called the length of the vector \(x)\) and is therefore independent of the chosen orientation on \(F'\) ; the sine and the tangent of \((x, y)\) change sign if one changes the orientation of \(F'\) Given two nonzero proportional vectors \(x, y\) of \(F\) , set \((x, y) = 0\) by convention.
% semantic-fragments: legacy-input-seam
\relax{}
